\documentclass[10pt,a4paper]{amsart}
\usepackage[margin=19mm]{geometry}
\usepackage{amsmath,amssymb,mathtools}
\usepackage[scr=rsfs,cal=euler]{mathalfa}
\usepackage{xfrac}
\usepackage[unicode,hidelinks,bookmarks=false]{hyperref}
\newcommand{\ie}{i.e.\ }
\newcommand{\cf}{cf.\ }
\newcommand{\resp}{respectively\ }
\newcommand{\Subsection}{\subsection}
\providecommand{\nobreakdash}{\nobreak\hbox{-}}
\setlength{\emergencystretch}{2em}
\multlinegap0pt
\usepackage{bm}
\usepackage{bbm}
\usepackage{dsfont}
\usepackage{xspace}
\usepackage{cancel}
\usepackage{mathtools}
\usepackage{amsthm}
\makeatletter
\@ifundefined{theorem}{\newtheorem{theorem}{Theorem}}{}
\@ifundefined{proposition}{\newtheorem{proposition}[theorem]{Proposition}}{}
\makeatother
\def\E{\mathbf{E}}
\def\H{{\mathcal H}}
\def\R{{\mathbb R}} % real numbers
\def\dist{\mathrm{dist}}
\def\ds{\displaystyle}
\def\grad{\nabla}
\newcommand{\pd}[2]{\frac{\partial #1}{\partial #2}}
\title[The Existence of Anomalous Dissipation over Bounded Interior]{The Existence of Anomalous Dissipation over Bounded Interior Domains}
\author{Ethan Dudley, Konstantina Trivisa}
\date{Source: arXiv:2501.00705v2 (CC BY 4.0)}
\begin{document}
\maketitle

\noindent\emph{Source and licence.} The Existence of Anomalous Dissipation over Bounded Interior Domains. Version arXiv:2501.00705v2 (CC BY 4.0), distributed under the Creative Commons Attribution 4.0 International licence, as recorded in the arXiv licence metadata for this exact version and shown on the abstract page https://arxiv.org/abs/2501.00705v2. Full rights notice: licenses/arxiv-2501.00705-CC-BY-4.0.txt.\par\smallskip

\noindent\textit{Editorial scope.} Counted unit: the Proposition whose source label is \texttt{proposition: choice of forcing} in the article (source0.tex lines 256--268, source label \texttt{proposition: choice of forcing}), delivered together with the article's own complete original proof of it (source0.tex lines 269--307, one \texttt{proof} environment of 39 lines). No mathematics is written, repaired or re-derived: statement and proof are the article's own text line for line. Required context retained uncounted: none. Every definition, display and label of those lines is retained unchanged. Omitted from this package and not redistributed: 1--255, 308--835. Nothing inside a retained excerpt is omitted or altered: every retained body is delivered exactly as the article's lines are written, so no omission receipt of any kind accompanies this package. Citations: the retained excerpts carry no citation command at all, so no bibliography entry of the article is transcribed and no reference is redistributed. Reference adaptation: none was needed and none was made; every reference inside the delivered unit resolves inside it, so no unresolved reference or citation is produced. Printed slips kept literal: no printed word, symbol, hypothesis or number of the counted statement or of its proof was altered. Layout and class: the article's own class is \texttt{article} with packages including amssymb, amsmath, amsthm, amsrefs, latexsym, color, esint, graphicx, subcaption; the wrapper is the lane's pinned \texttt{amsart} editorial wrapper at 10pt on A4, which restates the theorem-like environments the retained text needs and adds the article's own macro definitions that the retained text needs (\texttt{\textbackslash E}, \texttt{\textbackslash H}, \texttt{\textbackslash R}, \texttt{\textbackslash dist}, \texttt{\textbackslash ds}, \texttt{\textbackslash grad}, \texttt{\textbackslash pd}), with extra wrapper packages (bm, bbm, dsfont, xspace, cancel, mathtools). The delivered numbering is the unit's own. Assets: no retained span contains a figure environment, a graphics-inclusion command, a PDF-page inclusion, a table or a TikZ picture; no image file is copied into this package, the package declares no asset dependency and no image file is redistributed with it. Licence: version arXiv:2501.00705v2 of this article is distributed under the Creative Commons Attribution 4.0 International licence, as recorded in the arXiv licence metadata for this exact version and shown on the abstract page https://arxiv.org/abs/2501.00705v2. A full rights notice is delivered at licenses/arxiv-2501.00705-CC-BY-4.0.txt. Characters: every retained excerpt is pure ASCII, so the delivered engine, XeLaTeX with its default Latin Modern text font, reports no missing character.
\begin{proposition}\label{proposition: choice of forcing}
    Let $\delta \in (0,1)$, $R > 0$, $D = B(0,R)$, and $\{\mathbf{e}_j\}_{j=1}^3$ be the standard basis in $\R^3$. Let $\dist(x, A) = \ds\inf_{y \in A}\|x-y\|_{\ell^2}$ be the Euclidean distance to set the $A$. 
    Define
    \[
        g(x) := \frac{1}{\dist(x, \partial D)^{\delta/2}}\frac{-\sqrt{x_1^2 + x_2^2}\mathrm{e}_1 + x_3\mathrm{e}_2}{|x|} \quad x \in D.
    \]
    Then $g \in [L^2(0,T, L_\sigma^2(D))]^3$.
    Furthermore, for all $c,b,w > 0$ there exists a constant $C_\delta = C(c,b,\delta) >  0$ such that
    \begin{align*}
        \liminf_{\nu \to 0}\nu^{\delta/2}\E\int_0^T\int_{\partial D}\Big|\int_0^t\int_D \frac{c}{(\nu (t-s))^{3/2}}e^{-b\frac{|x-y|^2}{\nu (t-s)}}e^{-w\nu (t-s)}&g(y)dydW_s\Big|^2\H^2(dx)dt\\
        &\geq C_\delta T^{2-\delta/2} > 0.
    \end{align*}
\end{proposition}

\begin{proof}
    Notice that for $x \in B(0,R)$ then $\dist(x, \partial D) = R - |x|$ and $\widehat{\phi}(x) = \frac{-\sqrt{x_1^2 + x_2^2}\mathrm{e}_1 + x_3\mathrm{e}_2}{|x|}$ is the azimuthal unit vector. As such we can rewrite $g$ is spherical coordinates as 
    \[
        g(x) = \frac{1}{(R-r)^\delta}\widehat{\phi}
    \]
    where $|x| = r$. Then using spherical coordinates and the change of variables $R - r \mapsto r$ provides
    \begin{align*}
        \int_0^T\int_D|g|^2\;dxdt &= \int_0^T\int_0^R \int_{S^2}\frac{r^2}{(R - r)^\delta}\;dS(n)drdt\\
        &\leq 4\pi T\int_0^R \frac{(R-r)^2}{r^\delta}dr < \infty
    \end{align*}
    %    Now applying the triangle inequality to $g$ provides us with the uniform bound $g(x) \leq \frac{\sqrt{2}}{\sqrt{R}}$ for all $x \in D$
    and by applying the definition of divergence in spherical coordinates we can show that $g$ is divergence free:
    \[
        \grad \cdot g = \frac{1}{r\sin\theta}\pd{}{\phi}\Big(\frac{1}{(R - r)^{\delta/2}}\Big)= 0. 
    \]
    %As $L^2(D) \subset L^\infty(D)$ this implies $g \in L_\sigma^2(D)$. 
    
    Let $\chi_D(x)$ be the indicator function for the set $D$ and let $x \in \partial D$. Consider the change of variables $\xi = \frac{x-y}{\sqrt{\nu t}}$. Note the Triangle inequality implies that for all $h \in \R^3$
    \[
        |x-h| \leq |x| + |h| = R + |h| \Rightarrow |g(x-h)|^2 \geq \frac{1}{|h|^\delta}.
    \]
    Thus by applying Ito's isometry and our lower bound on $g(x-h)$ we obtain
    \begin{align*}
        S_\nu(t) &:= \nu^{\delta/2}\E\big|\int_0^t\int_D \frac{c}{(\nu (t-s))^{3/2}}e^{-b\frac{|x-y|^2}{\nu (t-s)}}e^{w\nu (t-s)}g(y)dydW_s\big|^2\\
        &= c^2\nu^{\delta/2} \E\big|\int_0^t\int_{\R^3} e^{-b|\xi|^2}e^{-w\nu (t-s)}g(x-\xi\sqrt{\nu (t-s)})\chi_D(x-\xi\sqrt{\nu (t-s)})\;d\xi dW_s\big|^2\\
        &=  c^2\nu^{\delta/2} \int_0^t\int_{\R^3} |e^{-b|\xi|^2}e^{-w\nu (t-s)}g(x-\xi\sqrt{\nu (t-s)})|^2\chi_D(x-\xi\sqrt{\nu (t-s)})\;d\xi ds\\
        &\geq c^2\nu^{\delta/2}\int_0^t\int_{\R^3} \frac{e^{-2b|\xi|^2}e^{-2w\nu (t-s)}}{|\xi|^\delta (\nu (t-s))^{\delta/2}}\chi_D(x-\xi\sqrt{\nu (t-s)})\;d\xi ds.
    \end{align*}
    Since $D = B(0, R)$ is convex, the radial lines $x - \xi\sqrt{\nu (t-s)}$ remain inside $D$ for $\nu$ sufficiently small whenever the angle between the vectors $x$ and $\xi$ is between $\frac{\pi}{2}$ and $\frac{3\pi}{2}$ radians. In other words if $\langle \cdot, \cdot \rangle_2$ is the $\ell^2$ inner product in $\R^3$ (i.e. the standard dot product in $\R^3$) then $x - \xi\sqrt{\nu (t-s)} \in D$ for some $\nu$ sufficiently small whenever $\langle x, \xi\rangle_2 < 0$. As such, the indicator functions $\chi_D(x - \xi\sqrt{\nu (t-s)}) \to \chi_{\langle \xi, x\rangle_2 < 0}$ in the sense of distributions as $\nu \to 0$ and for any fixed $x \in \partial D$, in the inviscid limit, the indicator function is supported over exactly half of the unit sphere (in $\xi$) relative to $x$.
    Hence by applying Fatou's Lemma and writing the $\xi$ integral in spherical coordinates, we obtain
    \begin{align*}
        \liminf_{\nu \to 0} \int_0^T\int_{\partial D} &S_\nu(t)\;\H^2(dx) dt\\
        &\geq \int_0^T\int_{\partial D} \liminf_{\nu \to 0}S_\nu(t)\;\H^2(dx)dt\\
        &= c^2 \int_0^T\int_0^t (t-s)^{-\delta/2}dsdt\int_{S^2} \int_{\{\xi \in \R^3 \mid \langle \xi, x\rangle_2 < 0\}} \frac{e^{-2b|\xi|^2}}{|\xi|^\delta}\;d\xi dS(x)\\
        &= \frac{c^2T^{2-\delta/2}}{(1-\delta/2)(2-\delta/2)}\int_{S^2}\int_{\{y \in S^2 \mid \langle y, x\rangle_2 < 0\}}\int_0^\infty e^{-2br^2}r^{2-\delta}\;drdS(y)dS(x)\\
        &= \frac{2\pi^2c^2T^{2-\delta/2}}{(2b)^{(3-\delta)/2}(1-\delta/2)(2-\delta/2)}\Gamma(3-\delta)\\
        &= C_\delta T^{2-\delta/2} > 0.
    \end{align*}
\end{proof}

\end{document}
